\section*{Chapter \ref{ch:pl:messaging-and-event-driven-architecture} --- Messaging and Event-Driven Architecture}

\begin{solution}{exo:pl:messaging-and-event-driven-architecture:1}
Order is kept only within a partition. Keying by account puts all of an account's fills in one partition, so the \omterm{def:pl:position-and-pnl-services:service}{position service} applies
them in the order they happened; round-robin would spread them over partitions read at different speeds.
\end{solution}

\begin{solution}{exo:pl:messaging-and-event-driven-architecture:2}
It resumes from the last committed offset, the start of the batch, and applies all 100 messages: the 30 already applied are applied twice.
\end{solution}

\begin{solution}{exo:pl:messaging-and-event-driven-architecture:3}
Round-robin: three, three and two partitions (0, 3, 6; 1, 4, 7; 2, 5). A fourth member triggers a rebalance to two each; each partition's new
owner resumes from the group's committed offset, so anything processed but not committed by the old owner is processed again.
\end{solution}

\begin{solution}{exo:pl:messaging-and-event-driven-architecture:4}
A crash falls uniformly within a batch of 100, on average half-way. Committing first has already moved the offset past the batch, so the
unprocessed half is lost; committing after has not, so the processed half is repeated. The observed 50\,323 and 48\,652 per 1\,000 crashes
are that half-batch.
\end{solution}

\begin{solution}{exo:pl:messaging-and-event-driven-architecture:5}
In the snapshot, written atomically with the state it describes: on restart it loads the snapshot and resumes from the offset stored in it,
replaying everything after. An offset committed elsewhere would not match the state.
\end{solution}

\begin{solution}{exo:pl:messaging-and-event-driven-architecture:6}
A relay that publishes an event and dies before marking the row sent publishes it again after the restart. The idempotent producer (the log
ignores a known outbox number) or the consumer (deduplicating on the event's key) removes it.
\end{solution}

\begin{solution}{exo:pl:messaging-and-event-driven-architecture:7}
Store the set of applied fill identifiers with the positions and skip any identifier already present; with the set and the positions updated
together, no fill is lost (the offset is still committed after) and none counts twice: zero and zero, as in the atomic mode.
\end{solution}

\begin{solution}{exo:pl:messaging-and-event-driven-architecture:8}
A messaging product's exactly-once covers what happens inside it --- deduplicated writes to its log, transactions across its partitions ---
not the effect of a message on another system. The position database is another system: its updates and the consumer's offset must be made
atomic, or its processing idempotent, by the consumer.
\end{solution}

\begin{solution}{pb:pl:messaging-and-event-driven-architecture:1}
\begin{enumerate}
\item Knowledge, time and rate: publishers do not know consumers, consumers can be down and catch up, a slow consumer does not slow the
publisher. It costs a visible call graph, and failures move to delivery.
\item A broker when storage, replay and changing consumers matter; brokerless (multicast, shared memory, a library such as Aeron) when latency
matters and the endpoints are known.
\item By its offset in its partition; the partition from a stable hash of the key.
\item The partitions of a topic, one member each, and one committed offset per partition; a rebalance reassigns partitions when members change.
\item Whole segments of old messages; a new subscriber that needs them must start from a snapshot that records the offset it covers.
\item Processing and committing the offset; a crash can fall inside the processing or between the two.
\item 967 in the commit-first mode, 1\,000 in the others: committing first never repeats messages, so a day has fewer steps and some scheduled
crash points fall after its end.
\item Commit first: 50\,323 lost, none duplicated; commit after: none lost, 48\,652 duplicated; atomic: none of either.
\item Up to 5\,100 shares (median 3\,500) committing first, up to 6\,500 (median 3\,800) committing after, zero atomic.
\item The positions and the next offset are written in one transaction; a crash rolls back both, and the consumer resumes where the state
ends.
\item They are two systems, the database and the log, with no transaction spanning both.
\item Write then publish: 22 events lost; publish then write: 22 phantom events.
\item Repeats (22 without idempotence keys), removed by the log's idempotent producer or by the consumer.
\item State stored as the sequence of events that changed it; the log must retain every event or a snapshot with its offset.
\item The messages behind a poison message in its partition.
\item \textbf{Named result.} Per 1\,000 crashes, committing before processing (at most once) loses 50\,323 fills and duplicates none;
committing after (at least once) loses none and duplicates 48\,652; storing positions and offset in one transaction loses and duplicates
none. The worst end-of-day position error over a hundred days is 5\,100, 6\,500 and 0 shares.
\item The atomic (or idempotent) mode for positions; a dashboard can commit first, since it only needs the latest state.
\item Its own log: deduplicated writes and transactions across its partitions, not the consumer's database.
\item Reconcile the positions with the drop copy continuously, per account and stock, and alert on any difference.
\item From the consumer, which makes its effect idempotent or atomic with its offset.
\end{enumerate}
\end{solution}

\begin{solution}{iq:pl:messaging-and-event-driven-architecture:1}
At most once: the position is saved before processing, so a crash loses messages. At least once: saved after, so a crash repeats them.
Exactly once: an effect, not a transport property, obtained by idempotent processing or by storing state and position atomically.
\iqlookfor{Where the offset is saved.}
\end{solution}

\begin{solution}{iq:pl:messaging-and-event-driven-architecture:2}
Commit the offset in the same database transaction as the positions, or store applied fill identifiers with them and skip repeats; never
rely on the transport alone.
\iqlookfor{Atomic offset or idempotence.}
\end{solution}

\begin{solution}{iq:pl:messaging-and-event-driven-architecture:3}
A \omterm{def:pl:messaging-and-event-driven-architecture:outbox}{transactional outbox}: the event is inserted in the same transaction as the row; a relay publishes unsent events in order and marks them,
with the outbox number as an idempotence key so that the repeats after a relay crash are dropped.
\iqlookfor{Outbox plus idempotence key.}
\end{solution}

\begin{solution}{iq:pl:messaging-and-event-driven-architecture:4}
Key the messages by account, so that one account's messages go to one partition, and let each partition be read by one consumer of the group
at a time.
\iqlookfor{Key to partition; one reader per partition.}
\end{solution}

\begin{solution}{iq:pl:messaging-and-event-driven-architecture:5}
A \omterm{def:pl:messaging-and-event-driven-architecture:log}{partitioned log} with topics per event type, keyed by account or trade; producers through outboxes; consumers idempotent or atomic with
their offsets; retention with snapshots for replay; dead letters monitored; consumer lag and \omterm{def:pl:post-trade-systems:recon}{reconciliations} against the drop copy as the
health measures.
\iqlookfor{Log, outbox, idempotent consumers, \omterm{def:pl:post-trade-systems:recon}{reconciliation}.}
\end{solution}
